% ==========================================
% BAB 2: DASAR TEORI
% ==========================================
\chapter{Dasar Teori}

\section{Version Control System (VCS) dan Git}
\textit{Version Control System} (VCS) adalah sistem yang merekam perubahan pada sebuah berkas atau sekumpulan berkas dari waktu ke waktu sehingga versi spesifik dapat dipanggil kembali di kemudian hari. Git merupakan salah satu VCS terdistribusi yang paling populer, dirancang untuk menangani proyek dari skala kecil hingga sangat besar dengan cepat dan efisien \cite{gitDoc}. Git memungkinkan kolaborasi tim yang lancar melalui fitur \textit{branching} (percabangan) dan \textit{merging} (penggabungan), di mana setiap pengembang dapat bekerja pada satu fitur secara terisolasi tanpa mengganggu basis kode utama.

\section{GitHub dan GitHub CLI}
GitHub adalah platform layanan \textit{hosting} repositori Git yang menyediakan antarmuka grafis berbasis web serta fitur pendukung kolaborasi seperti manajemen tugas dan integrasi berkelanjutan. GitHub memfasilitasi alur kerja pengembangan perangkat lunak modern melalui fitur \textit{Pull Request} (PR) dan \textit{Branch Protection Rules} untuk memastikan kualitas kode diulas terlebih dahulu sebelum digabungkan ke \textit{branch} utama \cite{githubDoc}. GitHub CLI (\texttt{gh}) adalah antarmuka baris perintah resmi untuk ekosistem GitHub yang memungkinkan pengembang untuk mengelola repositori, proses autentikasi, dan PR secara langsung dari terminal.

\section{Continuous Integration (CI) dan GitHub Actions}
\textit{Continuous Integration} (CI) adalah praktik rekayasa perangkat lunak di mana anggota tim secara teratur dan sering menggabungkan pekerjaan mereka ke dalam repositori terpusat, yang kemudian akan diverifikasi oleh proses \textit{build} dan pengujian otomatis \cite{cicdConcept}. Dalam ekosistem GitHub, praktik ini diimplementasikan menggunakan layanan \textbf{GitHub Actions}. Layanan ini memungkinkan pengembang untuk merancang alur kerja (\textit{workflow}) otomatis dalam sebuah berkas YAML guna melakukan kompilasi, pengujian logika (\textit{testing}), hingga standardisasi penulisan kode (\textit{linting}) setiap kali terjadi perubahan kode baru pada repositori.

\section{Framework Laravel}
Laravel adalah kerangka kerja (\textit{framework}) aplikasi web berbasis PHP yang mengimplementasikan arsitektur \textit{Model-View-Controller} (MVC). Laravel dirancang untuk membuat proses pengembangan web menjadi lebih cepat, aman, dan terstruktur dengan menyediakan fitur bawaan esensial seperti perutean (\textit{routing}), manajemen basis data (Eloquent ORM), dan modul pengujian terintegrasi menggunakan PHPUnit \cite{laravelDoc}. Dalam alur kerja otomatisasi modern, Laravel sering dipadukan dengan alat penjaga kualitas kode independen seperti Laravel Pint yang berbasis PHP-CS-Fixer untuk menjaga konsistensi format penulisan kode.

\chapter{Hasil dan Pembahasan}

\section{Hasil Praktikum}

\subsection{Preparation}

\subsubsection{Setup dan Autentikasi}
\begin{figure}[H]
    \centering
    \includegraphics[width=0.6\textwidth]{"install GitHub cli.png"}
    \caption{Install GitHub CLI}
    \label{fig:install_github_cli}
\end{figure}
Penjelasan Gambar:
\begin{itemize}
    \item Proses instalasi GitHub CLI pada sistem operasi Windows menggunakan perintah \texttt{winget}.
    \item Menampilkan proses pengunduhan dan verifikasi paket instalasi GitHub CLI hingga berhasil dipasang.
\end{itemize}

\begin{figure}[H]
    \centering
    \includegraphics[width=0.6\textwidth]{"gh auth login.png"}
    \caption{GH Auth Login}
    \label{fig:gh_auth_login}
\end{figure}
Penjelasan Gambar:
\begin{itemize}
    \item Melakukan pengecekan versi dan autentikasi akun GitHub menggunakan perintah \texttt{gh auth login}.
    \item Menunjukkan proses pemilihan protokol HTTPS dan autentikasi via peramban web dengan kode akses satu waktu (\textit{one-time code}).
\end{itemize}

\subsubsection{Manajemen Repositori dan Proteksi}
\begin{figure}[H]
    \centering
    \includegraphics[width=0.6\textwidth]{"make repo.png"}
    \caption{Make Repo}
    \label{fig:make_repo}
\end{figure}
Penjelasan Gambar:
\begin{itemize}
    \item Tampilan antarmuka repositori GitHub \texttt{evolusi-pl-24-542159-SV-25009} yang baru saja dibuat.
    \item Menampilkan struktur awal repositori yang berisi berkas \texttt{README.md} pada \textit{branch} \texttt{main}.
\end{itemize}

\begin{figure}[H]
    \centering
    \includegraphics[width=0.6\textwidth]{"github branch protection added to dev and main.png"}
    \caption{GitHub Branch Protection Added to Dev and Main}
    \label{fig:branch_protection}
\end{figure}
Penjelasan Gambar:
\begin{itemize}
    \item Pengaturan aturan proteksi \textit{branch} pada \textit{branch} utama (\texttt{main}) dan pengembangan (\texttt{dev}) di repositori GitHub.
    \item Memastikan setiap perubahan yang masuk harus melalui mekanisme Pull Request dan tidak bisa di-\textit{push} secara langsung.
\end{itemize}

\begin{figure}[H]
    \centering
    \includegraphics[width=0.6\textwidth]{"add identity.png"}
    \caption{Add Identity}
    \label{fig:add_identity}
\end{figure}
Penjelasan Gambar:
\begin{itemize}
    \item Penambahan konfigurasi identitas global Git menggunakan nama pengguna \texttt{Avin1731} dan email \texttt{les.017ciers@gmail.com}.
    \item Konfigurasi ini memastikan setiap \textit{commit} yang dilakukan terekam dengan identitas pengembang yang valid di repositori lokal.
\end{itemize}


\clearpage
\subsection{Penjelasan Singkat Perubahan Tiap Pull Request}

\subsubsection{Log Perubahan Fitur ke Branch Dev}
\begin{figure}[H]
    \centering
    \includegraphics[width=0.6\textwidth]{"log changes pr 1 feat to dev.png"}
    \caption{Log Changes PR 1 Feat to Dev}
    \label{fig:log_pr_1}
\end{figure}
Penjelasan Gambar:
\begin{itemize}
    \item Riwayat \textit{commit} pada Pull Request \#1 dari \textit{branch} \texttt{feature/splash-and-landing} ke \texttt{dev}.
    \item Menampilkan inisialisasi awal proyek Laravel 11, pembuatan \textit{splash screen} dengan animasi logo, serta modularisasi rute web publik dan admin.
\end{itemize}

\begin{figure}[H]
    \centering
    \includegraphics[width=0.6\textwidth]{"log changes pr 2 feat to dev.png"}
    \caption{Log Changes PR 2 Feat to Dev}
    \label{fig:log_pr_2}
\end{figure}
Penjelasan Gambar:
\begin{itemize}
    \item Riwayat \textit{commit} pada Pull Request \#2 dari \textit{branch} \texttt{feature/public-pages} ke \texttt{dev}.
    \item Menunjukkan penambahan halaman profil perusahaan publik, komponen navigasi, dan antarmuka galeri interaktif.
\end{itemize}

\begin{figure}[H]
    \centering
    \includegraphics[width=0.6\textwidth]{"log changes pr 3 feat to dev.png"}
    \caption{Log Changes PR 3 Feat to Dev}
    \label{fig:log_pr_3}
\end{figure}
Penjelasan Gambar:
\begin{itemize}
    \item Riwayat \textit{commit} pada Pull Request \#3 dari \textit{branch} \texttt{refactor/sync-splash-branding} ke \texttt{dev}.
    \item Memuat penyatuan identitas merek (\textit{brand identity}) pada halaman \textit{splash} dan publik, serta pembaruan dokumen \texttt{README}.
\end{itemize}

\begin{figure}[H]
    \centering
    \includegraphics[width=0.6\textwidth]{"log changes pr 4 feat to dev.png"}
    \caption{Log Changes PR 4 Feat to Dev}
    \label{fig:log_pr_4}
\end{figure}
Penjelasan Gambar:
\begin{itemize}
    \item Riwayat \textit{commit} pada Pull Request \#4 dari \textit{branch} \texttt{chore/setup-ci-workflow} ke \texttt{dev}.
    \item Menampilkan penambahan alur kerja GitHub Actions untuk pengujian (\textit{linting} dan \textit{test suites}), penyesuaian versi Node.js ke v22, serta penambahan target rilis v1.0.
\end{itemize}

\subsubsection{Log Perubahan ke Branch Main}
\begin{figure}[H]
    \centering
    \includegraphics[width=0.6\textwidth]{"log changes pr 5 dev to main.png"}
    \caption{Log Changes PR 5 Dev to Main}
    \label{fig:log_pr_5}
\end{figure}
Penjelasan Gambar:
\begin{itemize}
    \item Log riwayat \textit{commit} komprehensif saat menggabungkan seluruh fitur dari \textit{branch} \texttt{dev} menuju \texttt{main}.
    \item Menampilkan hierarki penggabungan (\textit{merge}) dari Pull Request \#1 hingga \#4 yang disatukan menjadi rilis produksi pada \textit{branch} utama.
\end{itemize}
\clearpage
\subsection{Penjelasan Singkat Log Commit}

\subsubsection{Log Commit Fitur dan Integrasi}
\begin{figure}[H]
    \centering
    \includegraphics[width=0.6\textwidth]{"log commit pr 1 (feat to dev).png"}
    \caption{Log Commit PR 1 (Feat to Dev)}
    \label{fig:commit_pr_1}
\end{figure}
Penjelasan Gambar:
\begin{itemize}
    \item Pesan \textit{commit} yang merekam riwayat pengembangan pada Pull Request \#1 sebelum digabungkan ke \texttt{dev}.
    \item Terdiri dari dua \textit{commit} utama, yaitu pembuatan \textit{splash screen} dan modularisasi rute web.
\end{itemize}

\begin{figure}[H]
    \centering
    \includegraphics[width=0.6\textwidth]{"log commit pr 2 (feat to dev).png"}
    \caption{Log Commit PR 2 (Feat to Dev)}
    \label{fig:commit_pr_2}
\end{figure}
Penjelasan Gambar:
\begin{itemize}
    \item Daftar pesan \textit{commit} pada pengerjaan Pull Request \#2.
    \item Mendokumentasikan penambahan halaman profil perusahaan publik beserta antarmuka galeri interaktif.
\end{itemize}

\begin{figure}[H]
    \centering
    \includegraphics[width=0.6\textwidth]{"log commit pr 3 (feat to dev).png"}
    \caption{Log Commit PR 3 (Feat to Dev)}
    \label{fig:commit_pr_3}
\end{figure}
Penjelasan Gambar:
\begin{itemize}
    \item Riwayat \textit{commit} pada Pull Request \#3 yang difokuskan pada perbaikan dan penyelarasan antarmuka.
    \item Melacak dua \textit{commit} terkait pembaruan identitas merek dan dokumentasi proyek secara komprehensif.
\end{itemize}

\begin{figure}[H]
    \centering
    \includegraphics[width=0.6\textwidth]{"log commit pr 4 (feat to dev).png"}
    \caption{Log Commit PR 4 (Feat to Dev)}
    \label{fig:commit_pr_4}
\end{figure}
Penjelasan Gambar:
\begin{itemize}
    \item Catatan \textit{commit} pada Pull Request \#4 yang mengatur integrasi berkelanjutan (CI) sebelum digabungkan ke \texttt{dev}.
    \item Memuat histori konfigurasi alur kerja GitHub Actions, pembaruan dokumen rilis, dan peningkatan versi Node.js.
\end{itemize}

\subsubsection{Log Commit Rilis dan Pratinjau}
\begin{figure}[H]
    \centering
    \includegraphics[width=0.6\textwidth]{"log commit pr 5 (dev to main).png"}
    \caption{Log Commit PR 5 (Dev to Main)}
    \label{fig:commit_pr_5}
\end{figure}
Penjelasan Gambar:
\begin{itemize}
    \item Kumpulan 13 pesan \textit{commit} pada proses penggabungan akhir dari \textit{branch} \texttt{dev} ke \texttt{main} untuk rilis v1.0.0.
    \item Riwayat migrasi kode stabil yang menyatukan seluruh fitur dan perbaikan yang telah diverifikasi.
\end{itemize}

\begin{figure}[H]
    \centering
    \includegraphics[width=0.6\textwidth]{"preview after push.png"}
    \caption{Preview After Push}
    \label{fig:preview_after_push}
\end{figure}
Penjelasan Gambar:
\begin{itemize}
    \item Tampilan pratinjau halaman repositori GitHub setelah perintah \textit{push} berhasil dijalankan dari terminal lokal.
    \item Munculnya notifikasi aktivitas \textit{push} terbaru beserta tombol akses cepat untuk membandingkan kode dan membuat Pull Request baru.
\end{itemize}
\clearpage

\subsection{Penjelasan Singkat Pembuatan Pull Request}

\subsubsection{Pembuatan Pull Request Fitur}
\begin{figure}[H]
    \centering
    \includegraphics[width=0.6\textwidth]{"create pr 1 feat to dev.png"}
    \caption{Create PR 1 Feat to Dev}
    \label{fig:create_pr_1}
\end{figure}
Penjelasan Gambar:
\begin{itemize}
    \item Menggunakan GitHub CLI (\texttt{gh pr create}) untuk membuat Pull Request pertama dari \textit{branch} \texttt{feature/splash-and-landing} ke \texttt{dev}.
    \item Menambahkan judul dan deskripsi terkait penambahan \textit{splash screen} interaktif dan modularisasi rute.
\end{itemize}

\begin{figure}[H]
    \centering
    \includegraphics[width=0.6\textwidth]{"create pr 2 feat to dev.png"}
    \caption{Create PR 2 Feat to Dev}
    \label{fig:create_pr_2}
\end{figure}
Penjelasan Gambar:
\begin{itemize}
    \item Eksekusi perintah CLI untuk pembuatan Pull Request kedua dari \textit{branch} \texttt{feature/public-pages} menuju \texttt{dev}.
    \item Menyertakan deskripsi penambahan berbagai halaman profil publik dan galeri interaktif.
\end{itemize}

\begin{figure}[H]
    \centering
    \includegraphics[width=0.6\textwidth]{"create pr 3 feat to dev.png"}
    \caption{Create PR 3 Feat to Dev}
    \label{fig:create_pr_3}
\end{figure}
Penjelasan Gambar:
\begin{itemize}
    \item Inisiasi Pull Request ketiga melalui CLI dari \textit{branch} \texttt{refactor/sync-splash-branding} ke \texttt{dev}.
    \item Berfokus pada pengisian detail penyatuan identitas merek dan pembaruan panduan pengaturan proyek.
\end{itemize}

\begin{figure}[H]
    \centering
    \includegraphics[width=0.6\textwidth]{"create pr 4 feat to dev.png"}
    \caption{Create PR 4 Feat to Dev}
    \label{fig:create_pr_4}
\end{figure}
Penjelasan Gambar:
\begin{itemize}
    \item Pembuatan Pull Request keempat menuju \textit{branch} \texttt{dev} menggunakan antarmuka CLI.
    \item Pengisian parameter judul dan deskripsi terkait konfigurasi GitHub Actions dan pengujian otomatis.
\end{itemize}

\subsubsection{Finalisasi dan Penutupan Pull Request}
\begin{figure}[H]
    \centering
    \includegraphics[width=0.6\textwidth]{"create pr 5 dev to main.png"}
    \caption{Create PR 5 Dev to Main}
    \label{fig:create_pr_5}
\end{figure}
Penjelasan Gambar:
\begin{itemize}
    \item Pembuatan Pull Request final dari \textit{branch} \texttt{dev} ke \texttt{main} menggunakan perintah GitHub CLI.
    \item Ditujukan sebagai tahap akhir untuk rilis resmi aplikasi versi 1.0.0 ke lingkungan produksi.
\end{itemize}

\begin{figure}[H]
    \centering
    \includegraphics[width=0.6\textwidth]{"pr list.png"}
    \caption{PR List}
    \label{fig:pr_list}
\end{figure}
Penjelasan Gambar:
\begin{itemize}
    \item Penggunaan perintah \texttt{gh pr list} pada terminal untuk melihat daftar Pull Request yang sedang aktif atau terbuka.
    \item Pemantauan status pengerjaan yang memperlihatkan Pull Request fitur pertama yang baru saja dibuat.
\end{itemize}

\begin{figure}[H]
    \centering
    \includegraphics[width=0.6\textwidth]{"menerima pr 1 feat ke dev.png"}
    \caption{Menerima PR 1 Feat ke Dev}
    \label{fig:merge_pr_1}
\end{figure}
Penjelasan Gambar:
\begin{itemize}
    \item Antarmuka web GitHub yang menunjukkan status Pull Request pertama siap untuk digabungkan (\textit{Able to merge}).
    \item Proses pratinjau dan persetujuan penggabungan kode fitur ke dalam \textit{branch} pengembangan \texttt{dev}.
\end{itemize}

\begin{figure}[H]
    \centering
    \includegraphics[width=0.6\textwidth]{"menerima pr dev ke main.png"}
    \caption{Menerima PR Dev ke Main}
    \label{fig:merge_dev_main}
\end{figure}
Penjelasan Gambar:
\begin{itemize}
    \item Tampilan halaman Pull Request rilis v1.0.0 di GitHub yang telah lolos verifikasi dan siap disetujui (\textit{Ready to merge}).
    \item Mengeksekusi penggabungan akumulasi 13 \textit{commit} dari seluruh fitur \texttt{dev} menuju \texttt{main}.
\end{itemize}

\begin{figure}[H]
    \centering
    \includegraphics[width=0.6\textwidth]{"pr closed.png"}
    \caption{PR Closed}
    \label{fig:pr_closed}
\end{figure}
Penjelasan Gambar:
\begin{itemize}
    \item Daftar status kelima Pull Request yang telah berhasil ditutup dan digabungkan (\textit{merged}) pada repositori.
    \item Menandakan bahwa seluruh siklus integrasi fitur dan persiapan rilis v1.0.0 telah sukses diselesaikan.
\end{itemize}

\clearpage
\subsection{Penjelasan Singkat ci.yml}

\subsubsection{Konfigurasi Awal dan Kendala GitHub Actions}

Berkas konfigurasi \texttt{ci.yml} berisi 2 \textit{job} utama, yaitu Laravel Pint dan PHPUnit untuk otomatisasi pengujian kode. Berikut adalah rincian cara kerja dari kedua \textit{job} tersebut dalam alur kerja (CI):

\begin{enumerate}
    \item \textbf{Job 1: Code Quality \& Style (Laravel Pint (\texttt{lint-code}))}
    \begin{itemize}
        \item \textbf{Tujuan:} Memeriksa kerapian, konsistensi, dan standar penulisan kode PHP pada proyek agar sesuai dengan aturan yang berlaku (Laravel Pint berbasis PHP-CS-Fixer).
        \item \textbf{Cara Kerja:} \textit{Workflow} akan menyiapkan lingkungan PHP versi 8.3 dengan ekstensi pendukung, menginstal dependensi melalui Composer, lalu menjalankan perintah pengujian gaya kode \texttt{./vendor/bin/pint --test}. Jika ditemukan format kode yang tidak sesuai, \textit{job} ini akan langsung menggagalkan proses CI untuk menjaga kebersihan kode.
    \end{itemize}
    
    \item \textbf{Job 2: Automated Unit \& Feature Tests (PHPUnit (\texttt{run-tests}))}
    \begin{itemize}
        \item \textbf{Tujuan:} Menjalankan rangkaian pengujian otomatis (\textit{unit and feature test suite}) untuk memastikan seluruh logika bisnis dan fungsionalitas aplikasi berjalan dengan benar.
        \item \textbf{Cara Kerja:} \textit{Workflow} menyiapkan lingkungan PHP 8.3 (lengkap dengan ekstensi SQLite), menyiapkan Node.js beserta manajer paket PNPM untuk membangun aset frontend melalui Vite, menginstal dependensi backend dan frontend, menyalin konfigurasi lingkungan (\texttt{.env}), lalu mengeksekusi pengujian menggunakan basis data SQLite in-memory melalui perintah \texttt{php artisan test}.
    \end{itemize}
\end{enumerate}

Sempat terjadi galat (\textit{error}) pada implementasi awal seperti pada gambar di bawah ini:

\begin{figure}[H]
    \centering
    \begin{minipage}[b]{0.48\textwidth}
        \centering
        \includegraphics[width=\textwidth]{"code b4 1.png"}
    \end{minipage}
    \hfill
    \begin{minipage}[b]{0.48\textwidth}
        \centering
        \includegraphics[width=\textwidth]{"code b4 2.png"}
    \end{minipage}
    \caption{Konfigurasi Awal Berkas ci.yml (Kiri: Bagian 1, Kanan: Bagian 2)}
    \label{fig:ci_yml_b4}
\end{figure}

Penjelasan Gambar:
\begin{itemize}
    \item Tampilan kode bagian awal dan akhir dari \texttt{ci.yml} sebelum dilakukan penyesuaian versi Node.js.
    \item Menampilkan struktur \textit{workflow} secara utuh untuk \textit{job} Laravel Pint dan PHPUnit.
\end{itemize}

\begin{figure}[H]
    \centering
    \includegraphics[width=0.6\textwidth]{"github action 1 still error.png"}
    \caption{GitHub Action 1 Still Error}
    \label{fig:ga_error_1}
\end{figure}
Penjelasan Gambar:
\begin{itemize}
    \item Sempat terjadi galat (\textit{error}) pada eksekusi awal GitHub Actions.
    \item Indikasi kegagalan pada proses pengujian otomatis.
\end{itemize}

\begin{figure}[H]
    \centering
    \includegraphics[width=0.6\textwidth]{"github action 1 still error .2.png"}
    \caption{GitHub Action 1 Still Error .2}
    \label{fig:ga_error_2}
\end{figure}
Penjelasan Gambar:
\begin{itemize}
    \item Lanjutan kendala atau galat yang sama pada iterasi pengujian GitHub Actions berikutnya.
    \item Analisis pesan kesalahan pada konsol tindakan.
\end{itemize}

\begin{figure}[H]
    \centering
    \includegraphics[width=0.6\textwidth]{"log error github action.png"}
    \caption{Log Error GitHub Action}
    \label{fig:log_error_ga}
\end{figure}
Penjelasan Gambar:
\begin{itemize}
    \item Rincian log kesalahan (\textit{error log}) pada GitHub Actions yang menunjukkan baris atau proses yang bermasalah.
    \item Referensi utama untuk melakukan perbaikan kode konfigurasi maupun skrip pengujian.
\end{itemize}

\subsubsection{Hasil Perbaikan dan Pengujian Berhasil}

Lalu setelah diperbaiki, berikut adalah kode dan hasil eksekusinya:

\begin{figure}[H]
    \centering
    \includegraphics[width=0.6\textwidth]{"code after.png"}
    \caption{Code After}
    \label{fig:code_after}
\end{figure}
Penjelasan Gambar:
\begin{itemize}
    \item Perbaikan pada kode konfigurasi dengan memperbarui versi Node.js menjadi \texttt{22} agar kompatibel dengan dependensi \texttt{pnpm}.
    \item Hasil perubahan (\textit{diff}) dari berkas \texttt{ci.yml}.
\end{itemize}

\begin{figure}[H]
    \centering
    \includegraphics[width=0.6\textwidth]{"github action success.png"}
    \caption{GitHub Action Success}
    \label{fig:ga_success}
\end{figure}
Penjelasan Gambar:
\begin{itemize}
    \item Tampilan status GitHub Actions yang menunjukkan keberhasilan setelah perbaikan diterapkan.
    \item Seluruh rangkaian pengujian lolos tanpa kendala.
\end{itemize}

\begin{figure}[H]
    \centering
    \includegraphics[width=0.6\textwidth]{"laravel pint success.png"}
    \caption{Laravel Pint Success}
    \label{fig:pint_success}
\end{figure}
Penjelasan Gambar:
\begin{itemize}
    \item Hasil eksekusi \textit{job} Laravel Pint yang menunjukkan bahwa format dan gaya penulisan kode PHP telah sesuai standar.
    \item Tidak ditemukan pelanggaran gaya kode (\textit{code style}).
\end{itemize}

\begin{figure}[H]
    \centering
    \includegraphics[width=0.6\textwidth]{"php unit success.png"}
    \caption{PHP Unit Success}
    \label{fig:phpunit_success}
\end{figure}
Penjelasan Gambar:
\begin{itemize}
    \item Laporan pengujian unit menggunakan PHPUnit yang berhasil melewati seluruh skenario pengujian (\textit{test suite}).
    \item Memastikan fungsi logika program berjalan dengan benar.
\end{itemize}

\begin{figure}[H]
    \centering
    \includegraphics[width=0.6\textwidth]{"log github action is running.png"}
    \caption{Log GitHub Action Is Running}
    \label{fig:ga_running}
\end{figure}
Penjelasan Gambar:
\begin{itemize}
    \item Catatan log \textit{real-time} yang menunjukkan proses GitHub Actions sedang berjalan secara aktif.
    \item Pemantauan tahapan eksekusi \textit{job} secara langsung pada sistem jarak jauh.
\end{itemize}

\clearpage
\section{Tautan Repositori Kode Sumber}
Seluruh \textit{source code} dari proyek aplikasi web berbasis Laravel dan konfigurasi alur kerja GitHub Actions pada praktikum ini dapat diakses secara publik melalui tautan repositori GitHub berikut:
\begin{itemize}
    \item \textbf{URL Repositori:} \url{https://github.com/Avin1731/evolusi-pl-24-542159-SV-25009}
\end{itemize}
\noindent\textbf{Catatan Tambahan terkait Status Repositori:}\\
Saat laporan ini disusun, repositori basis kode masih di-\textit{hosting} pada akun GitHub pribadi praktikan. Repositori tersebut saat ini belum dapat ditransfer atau diserahkan ke dalam GitHub Organisasi resmi milik kelas atau mata kuliah praktikum. Hal ini dikarenakan praktikan belum menerima undangan (\textit{invite}) untuk bergabung ke dalam organisasi tersebut, sehingga praktikan belum memiliki hak akses (\textit{permission}) administratif yang dibutuhkan untuk mengeksekusi proses transfer repositori.

\chapter{Kesimpulan}
Berdasarkan praktikum Konstruksi dan Evolusi Perangkat Lunak yang telah dilakukan, dapat ditarik beberapa kesimpulan sebagai berikut:
\begin{enumerate}
    \item Penggunaan Git terdistribusi bersama GitHub CLI (\texttt{gh}) sangat efektif dalam menyederhanakan dan mempercepat alur kerja kontrol versi. Melalui terminal, tahapan inisialisasi repositori, autentikasi sesi, hingga manajemen \textit{Pull Request} dapat dieksekusi tanpa perlu berpindah konteks ke aplikasi web.
    \item Penerapan strategi percabangan kode (\textit{branching}) yang memisahkan antara lingkungan \texttt{main}, \texttt{dev}, dan \texttt{feature/*}, yang dikombinasikan dengan aturan \textit{Branch Protection}, terbukti berhasil melindungi integritas basis kode. Mekanisme ini mencegah masuknya modifikasi kode secara sembarangan yang berpotensi merusak lingkungan produksi.
    \item Otomatisasi alur kerja menggunakan \textit{Continuous Integration} (CI) via GitHub Actions (\texttt{ci.yml}) dapat memastikan setiap integrasi kode baru divalidasi dengan ketat. Sistem secara otomatis menjalankan \textit{code styling} menggunakan Laravel Pint dan pengujian logika melalui PHPUnit tanpa intervensi manual, sehingga stabilitas aplikasi tetap terjaga.
    \item Penggabungan antara kerangka kerja PHP modern (Laravel), manajemen kode sumber kolaboratif (GitHub), dan jalur otomatisasi pengujian (CI/CD) adalah fondasi esensial untuk menyimulasikan ekosistem pengembangan rekayasa perangkat lunak di tingkat industri saat ini.
\end{enumerate}
% ==========================================
% DAFTAR PUSTAKA
% ==========================================
\renewcommand{\bibname}{Daftar Pustaka}
\begin{thebibliography}{9}
\bibitem{gitDoc} 
Chacon, S., \& Straub, B. (2014). \textit{Pro Git} (2nd ed.). Apress. Tersedia di: \url{https://git-scm.com/book/en/v2}
\bibitem{githubDoc} 
GitHub Docs. (2024). \textit{About collaborative development models}. Tersedia di: \url{https://docs.github.com/en/pull-requests/collaborating-with-pull-requests}
\bibitem{cicdConcept} 
Fowler, M. (2006). \textit{Continuous Integration}. Tersedia di: \url{https://martinfowler.com/articles/continuousIntegration.html}
\bibitem{laravelDoc} 
Otwell, T. (2024). \textit{Laravel Documentation: Release 11.x}. Tersedia di: \url{https://laravel.com/docs/11.x}
\end{thebibliography}